% Figure: an electrodialysis stack schematic. Alternating cation- and
% anion-exchange membranes between an anode and a cathode split a feed into
% alternating dilute and concentrate channels; ions migrate under the applied
% field toward their counter-electrode, are stopped by the membrane of the
% wrong charge, and so accumulate in the concentrate channels while the dilute
% channels are stripped. A cell pair is one CEM and one AEM plus their two
% channels; a stack has many cell pairs between one electrode pair.
\begin{figure}[htbp]
\centering
\begin{tikzpicture}[>=stealth, font=\small]
% electrodes
\fill[enggray!30] (0,0) rectangle (0.35,4);
\node[enggray, font=\scriptsize, rotate=90] at (0.175,2) {anode $+$};
\fill[enggray!30] (8.65,0) rectangle (9,4);
\node[enggray, font=\scriptsize, rotate=90] at (8.825,2) {cathode $-$};
% membranes: C = cation-exchange (blue), A = anion-exchange (red)
\foreach \x/\lab/\col in {1.5/C/engblue, 3/A/engred, 4.5/C/engblue, 6/A/engred, 7.5/C/engblue}{
  \draw[\col, very thick] (\x,0) -- (\x,4);
  \node[\col, font=\scriptsize] at (\x,4.25) {\lab};
}
% channel labels: between electrodes the channels alternate dilute / concentrate
\node[engblue, font=\scriptsize] at (0.95,-0.35) {conc};
\node[englime!70!black, font=\scriptsize] at (2.25,-0.35) {dilute};
\node[engblue, font=\scriptsize] at (3.75,-0.35) {conc};
\node[englime!70!black, font=\scriptsize] at (5.25,-0.35) {dilute};
\node[engblue, font=\scriptsize] at (6.75,-0.35) {conc};
\node[engblue, font=\scriptsize] at (8.2,-0.35) {conc};
% ion motion arrows: cations (+) move toward cathode (right), anions (-) left.
% show one cation and one anion in a dilute channel
\draw[->, engblue, thick] (5.0,2.6) -- (5.55,2.6);
\node[engblue, font=\scriptsize] at (5.05,2.9) {$+$};
\draw[->, engred, thick] (5.0,1.4) -- (4.55,1.4);
\node[engred, font=\scriptsize] at (5.05,1.1) {$-$};
% bracket for one cell pair
\draw[decorate, decoration={brace, amplitude=5pt}, enggray]
  (3,-0.75) -- (6,-0.75);
\node[enggray, font=\scriptsize] at (4.5,-1.15) {one cell pair};
\node[font=\scriptsize, anchor=west] at (0.2,4.7)
  {C: cation-exchange membrane \quad A: anion-exchange membrane};
\end{tikzpicture}
\caption[An electrodialysis stack of alternating ion-exchange
membranes]{An electrodialysis stack. Alternating cation-exchange (C) and
anion-exchange (A) membranes sit between one electrode pair. Under the applied
field cations migrate toward the cathode and anions toward the anode; each ion
passes the membrane that admits its own charge and is stopped by the membrane
of the opposite charge, so ions are pumped out of the dilute channels and
trapped in the neighbouring concentrate channels. One cation-exchange plus one
anion-exchange membrane and their two channels form a cell pair, and a
production stack holds many cell pairs in series so that a single electrode
current desalts every dilute channel at once. The transport is driven by the
electric potential rather than by a pressure or a concentration difference,
which makes electrodialysis efficient for brackish water where the salt to be
removed is a small fraction of the stream.}
\label{fig:vol04ch07:electrodialysis}
\end{figure}
